\chapter{Causal Forests}
\label{chap:causal-forest}

\section*{Outline}

The ATE summarizes an average effect, but policy questions often concern
heterogeneity: who benefits, who is harmed, and how does the effect vary with
observed characteristics?  This chapter introduces the conditional average
treatment effect (CATE), reviews why conventional nonparametric regression becomes
difficult in many dimensions, and explains the ideas behind honest causal trees and
causal forests.

\section{CATE and treatment-effect heterogeneity}
\label{sec:cf-cate}

For a binary treatment, the average treatment effect is
\[
  \tau=\E[Y_i(1)-Y_i(0)].
\]
It conceals heterogeneity.  The individual treatment effect
$Y_i(1)-Y_i(0)$ contains all such variation, but it cannot normally be observed and
does not by itself explain which characteristics generate the variation.

Let $W_i$ be a vector of pre-treatment characteristics.  The
\emph{conditional average treatment effect} is
\begin{equation}
  \tau(w)\equiv \operatorname{CATE}(w)
  =\E[Y_i(1)-Y_i(0)\mid W_i=w].
  \label{eq:cf-cate}
\end{equation}
This object organizes heterogeneity by observed characteristics.  If $W_i$ is a
gender indicator, for example, comparing $\tau(1)$ and $\tau(0)$ describes how the
effect differs between those groups.  The law of iterated expectations also gives
\begin{equation}
  \tau=\E[\tau(W_i)].
  \label{eq:cf-ate-from-cate}
\end{equation}
Thus the CATE function contains the information needed to recover the ATE.

\subsection{Unconfoundedness, overlap, and identification}
\label{sec:cf-identification}

Randomized treatment is not necessary if treatment is as good as random after
conditioning on observed covariates.  The key assumption is
\begin{equation}
  D_i\indep\bigl(Y_i(1),Y_i(0)\bigr)\mid W_i,
  \label{eq:cf-unconfoundedness}
\end{equation}
called \emph{unconfoundedness}, \emph{conditional exchangeability}, or
\emph{selection on observables}.  We also need overlap at the covariate value of
interest:
\begin{equation}
  0<\Pr(D_i=1\mid W_i=w)<1.
  \label{eq:cf-overlap}
\end{equation}
Without overlap, one of the two conditional outcome means cannot be learned from
the data at that $w$.

Complete randomization of $D_i$ independently of potential outcomes and $W_i$
implies \eqref{eq:cf-unconfoundedness}, but the converse need not hold.  Stratified
experiments are a useful example: treatment probabilities may differ by gender,
age, or site, while assignment remains random within each stratum.  In observational
data, \eqref{eq:cf-unconfoundedness} requires all common causes of treatment and the
potential outcomes to be included in $W_i$.  It is a substantive assumption, not
something that can be verified from the observed data alone.

For intuition, suppose $Y_i(d)=g_d(W_i,U_i)$.  A sufficient condition is that,
conditional on $W_i$, treatment is independent of the outcome-relevant unobservable
$U_i$.  Merely allowing treatment to remain related to such an $U_i$ after
conditioning on $W_i$ would generally violate unconfoundedness.

Under \eqref{eq:cf-unconfoundedness} and \eqref{eq:cf-overlap},
\begin{align}
  \tau(w)
  &=\E[Y_i(1)\mid W_i=w]-\E[Y_i(0)\mid W_i=w] \notag\\
  &=\E[Y_i(1)\mid D_i=1,W_i=w]
    -\E[Y_i(0)\mid D_i=0,W_i=w] \notag\\
  &=\E[Y_i\mid D_i=1,W_i=w]
    -\E[Y_i\mid D_i=0,W_i=w].
  \label{eq:cf-identification}
\end{align}
The second equality uses unconfoundedness and the last uses consistency.  Define
\[
  \mu_d(w)=\E[Y_i\mid D_i=d,W_i=w].
\]
Estimating the CATE is now a problem of estimating the two regression functions
$\mu_1$ and $\mu_0$, or estimating their contrast directly.

If $W$ is discrete, a cell mean may work.  If $W$ is continuous, the event $W=w$
usually has probability zero, so exact matching is unavailable.  We need a method
that borrows information across nearby or otherwise similar observations.

\section{Traditional nonparametric methods}
\label{sec:cf-nonparametrics}

There are two broad ways to estimate an unknown regression function.

\subsection{Local methods}

A local method estimates $\mu_d(w)$ using observations with $W_i$ near $w$.  With a
scalar $W$, a simple estimator averages observations for which
$|W_i-w|\leq h$.  The bandwidth $h$ controls a familiar trade-off:

\begin{itemize}
  \item a larger $h$ uses more observations and reduces variance, but introduces
  more approximation bias;
  \item a smaller $h$ focuses more tightly on $w$ and reduces bias, but increases
  variance.
\end{itemize}

If $W$ has dimension $p$ and each coordinate uses a bandwidth of order $h$, the
effective sample size is of order $nh^p$.  Consequently, a typical pointwise
standard error is of order
\begin{equation}
  \frac{1}{\sqrt{nh^p}},
  \label{eq:cf-local-se}
\end{equation}
not $\sqrt{nh^p}$.  As $p$ grows, it becomes difficult to find observations close
to $w$ in every coordinate.  This is the \emph{curse of dimensionality}.

\subsection{Global methods}

A global method approximates the entire function.  For known basis functions
$b_1,b_2,\ldots$, one may write
\[
  \mu_d(w)\approx \sum_{j=1}^{J}a_{dj}b_j(w).
\]
Polynomials, B-splines, and wavelets are common choices.  Increasing $J$ reduces
approximation bias but requires estimating more coefficients.  With several
covariates, the number of basis terms needed for a flexible approximation can grow
quickly, so the curse of dimensionality remains.

Many machine-learning methods fit within this local--global distinction.  Trees and
forests create adaptive neighborhoods and can be viewed as adaptive local methods.
Neural networks, studied later, construct a global approximation.

\section{Regression trees, causal trees, and causal forests}
\label{sec:cf-trees}

The central insight of a tree is that geometric closeness is sufficient but not
necessary for similar regression values.  Two distant covariate vectors may have
the same $\mu_d(w)$, while a function may change sharply over a small distance.  A
tree uses the outcomes to learn rectangular regions within which the target is
approximately constant.

\subsection{A regression-tree sketch}

Suppose $W=(W_1,W_2)'$ lies in the unit square.  A simplified recursive partitioning
algorithm proceeds as follows.

\begin{enumerate}
  \item For each covariate $W_j$ and each admissible threshold $t$, split the current
  node into $W_j\leq t$ and $W_j>t$.
  \item Evaluate how much the split reduces within-node squared error, equivalently
  how strongly the two child-node means differ after accounting for their sizes.
  \item Choose the best admissible split and repeat within each child node.
  \item Stop when a depth or minimum-leaf-size rule is reached.  Predict at $w$ by
  averaging outcomes in the terminal leaf containing $w$.
\end{enumerate}

The one-coordinate-at-a-time splits make computation manageable and give an
interpretable binary-tree representation.  In practice, candidate thresholds are
the finitely many cut points induced by the sample, and split criteria include
regularization to avoid tiny or severely imbalanced leaves.

\begin{figure}[htbp]
  \centering
  \includegraphics[width=0.94\textwidth]{illuTree.png}
  \caption{A generic partition, a recursive rectangular partition, its tree
  representation, and the resulting piecewise-constant fit.  Adapted from
  Friedman, Hastie, and Tibshirani, \emph{The Elements of Statistical Learning}
  (2009), p.~306.}
  \label{fig:cf-tree}
\end{figure}

In Figure~\ref{fig:cf-tree}, the top-left partition is not generally representable
by recursive one-coordinate splits.  The top-right partition is; the bottom-left
panel records the corresponding decision rules, and the bottom-right panel shows
the fitted step function.

\subsection{From a regression tree to a causal tree}
\label{sec:cf-causal-tree}

A causal tree targets the contrast in \eqref{eq:cf-identification}, not merely an
outcome mean.  In a simple description, each proposed child node contains a treated
mean and a control mean, and hence an estimated treatment effect.  A useful split is
one for which the treatment-effect estimates differ substantially across the two
children, subject to penalties for estimation noise and constraints ensuring that
both treatment arms are represented.

This description conveys the intuition, but production implementations use more
carefully regularized splitting objectives than the raw squared difference between
two noisy effect estimates.

\subsection{Honesty}
\label{sec:cf-honesty}

If the same outcomes select the partition and estimate effects inside the selected
leaves, the procedure can chase noise.  Conditional on the selected leaf, the
within-leaf outcomes are no longer an ordinary, exogenously selected sample.
\emph{Honesty} separates the two tasks:

\begin{enumerate}
  \item randomly divide a subsample into a splitting sample and an estimation
  sample;
  \item use only the splitting sample to construct the tree;
  \item hold the partition fixed and use only the estimation sample to estimate
  effects inside its leaves.
\end{enumerate}

Under suitable regularity conditions---including overlap, regular splits, growing
leaf information, and smoothness---honest tree and forest estimators can be shown
to be consistent and asymptotically normal.  These are asymptotic statements, not a
claim of exact finite-sample unbiasedness.  Pointwise nonparametric rates are
typically slower than the parametric $n^{-1/2}$ rate.

\subsection{Causal forests}

A single causal tree is unstable: a modest change in the data can alter an early
split and therefore the entire partition.  A causal forest averages many randomized
honest trees:

\begin{enumerate}
  \item draw a subsample of the observations;
  \item grow an honest causal tree, usually considering a random subset of
  covariates at each split;
  \item repeat for $B$ trees;
  \item average the $B$ tree predictions at the target $w$.
\end{enumerate}

Subsampling and random covariate selection make the trees less correlated; averaging
then stabilizes the prediction.  Across trees, an observation may be used for
splitting in some trees, effect estimation in others, and omitted from still others.

Three tuning ideas are especially important:

\begin{description}
  \item[Covariate subsampling.] Considering only a random subset of the $p$
  covariates at a split reduces computation and decorrelates the trees.  The
  \texttt{grf} package chooses sensible defaults that can be tuned.
  \item[Minimum leaf size.] Large leaves can smooth over real heterogeneity; tiny
  leaves have noisy treated and control means.
  \item[Balanced splits.] Each child needs enough observations, including enough
  treated and control observations, for stable effect estimation and overlap.
\end{description}

\subsection{Why ordinary linear regression is not enough}

Under complete randomization, the treatment coefficient in a short regression
identifies the ATE even if the outcome model is nonlinear.  CATE is different: it is
a function of $W$.  A linear interaction model such as
\[
  Y=\alpha+D\tau+W'\beta+D W'\delta+u
\]
restricts the CATE to $\tau+W'\delta$.  If the true CATE is nonlinear, this function
is misspecified.  Moreover, under unconfoundedness $D$ may be related to $W$, so
omitted nonlinear functions of $W$ can also confound treatment.  Causal forests
allow flexible nonlinearities and interactions, though they do not remove the need
for unconfoundedness and overlap.

\section{Implementation with \texttt{grf}}
\label{sec:cf-implementation}

The next simulation has $n=2{,}000$, ten covariates, and a treatment probability
that changes with $W_1$.  The outcome model is
\[
  Y=\max\{W_1,0\}D+W_2+\min\{W_3,0\}+\varepsilon,
\]
so the true CATE is $\tau(w)=\max\{w_1,0\}$.

\begin{rcode}
library(grf)
library(DiagrammeR)

set.seed(5280)
n <- 2000
p <- 10
W <- matrix(rnorm(n * p), nrow = n, ncol = p)
D <- rbinom(n, 1, 0.4 + 0.2 * (W[, 1] > 0))
Y <- pmax(W[, 1], 0) * D + W[, 2] + pmin(W[, 3], 0) + rnorm(n)

W_test <- matrix(0, nrow = 101, ncol = p)
W_test[, 1] <- seq(-2, 2, length.out = 101)

# A quick exploratory forest and one of its trees
forest_quick <- causal_forest(W, Y, D)
plot(get_tree(forest_quick, 1))

# A larger forest for final estimates and variance estimates
forest_full <- causal_forest(W, Y, D, num.trees = 4000)
pred <- predict(forest_full, W_test, estimate.variance = TRUE)
se <- sqrt(pred$variance.estimates)

plot(W_test[, 1], pred$predictions,
     ylim = range(pred$predictions - 1.96 * se,
                  pred$predictions + 1.96 * se, 0, 2),
     xlab = expression(W[1]), ylab = "CATE", type = "l")
lines(W_test[, 1], pred$predictions + 1.96 * se, lty = 2)
lines(W_test[, 1], pred$predictions - 1.96 * se, lty = 2)
lines(W_test[, 1], pmax(W_test[, 1], 0), col = 2)
legend("topleft", c("forest", "95% pointwise interval", "truth"),
       lty = c(1, 2, 1), col = c(1, 1, 2), bty = "n")
\end{rcode}

Four thousand trees reduce Monte Carlo noise in variance estimates but can be slow
in a browser.  The live companion uses fewer trees for a quick classroom
demonstration; the substantive workflow should use more trees and check that results
are stable as \texttt{num.trees} increases.

We can also revisit the survey-wording experiment from
Section~\ref{sec:rct-gss}.  Here the test grid fixes age at 30 and varies education.

\begin{rcode}
gss <- read.csv("data/welfare-small.csv")
Y <- as.numeric(gss$Y)
D <- as.numeric(gss$D)
W <- as.matrix(gss[, c("age", "educ")])

W_test <- cbind(age = 30,
                educ = seq(min(gss$educ), max(gss$educ), by = 1))

forest_gss <- causal_forest(W, Y, D, num.trees = 4000)
pred_gss <- predict(forest_gss, W_test, estimate.variance = TRUE)
se_gss <- sqrt(pred_gss$variance.estimates)

plot(W_test[, "educ"], pred_gss$predictions,
     ylim = range(pred_gss$predictions - 1.96 * se_gss,
                  pred_gss$predictions + 1.96 * se_gss),
     xlab = "Education", ylab = "Estimated CATE", type = "l")
lines(W_test[, "educ"], pred_gss$predictions + 1.96 * se_gss, lty = 2)
lines(W_test[, "educ"], pred_gss$predictions - 1.96 * se_gss, lty = 2)
\end{rcode}

The confidence limits above are pointwise.  They should not be interpreted as a
simultaneous confidence band for the whole CATE curve.  In an observational
application, predictive fit also cannot validate unconfoundedness.

\section{Summary}

\begin{itemize}
  \item The CATE organizes treatment-effect heterogeneity by pre-treatment
  characteristics, and averaging it recovers the ATE.
  \item Unconfoundedness and overlap identify the CATE as a contrast of two
  observable conditional means.
  \item Conventional nonparametric methods face the curse of dimensionality; a
  local estimator's standard error is typically of order $(nh^p)^{-1/2}$.
  \item Causal trees learn adaptive regions targeted to treatment-effect
  heterogeneity.  Honesty separates partition selection from effect estimation.
  \item Causal forests average many randomized honest trees, improving stability
  and permitting pointwise inference under regularity conditions.
  \item Flexible algorithms do not make causal assumptions disappear: credible
  unconfoundedness, overlap, SUTVA, and a clearly defined target population remain
  essential.
\end{itemize}
